\subsection*{Question 1}
	On range les valeurs de la série dans l'ordre croissant :

	\centerline{\begin{tblr}{colspec={*{7}{X[c,1cm]}},hline{1,2}={4}{solid},vline{4,5}={solid}}
		7&7&9&9&12&23&25\\
	\end{tblr}}

	Comme l'effectif total est impair (il y a 7 valeurs), la médiane est la valeur centrale (la quatrième), dans l'ordre croissant.

	La médiane de la série est 9.

\subsection*{Question 2}
	Le 4 qui est le premier chiffre non nul est le quatrième chiffre après la virgule, le chiffre des dix-millièmes.

	On a: \quad $\np{0,000457} = 4,57 \times \np{0,0001} = 4,57 \times 10^{-4}$.

	La notation scientifique de $\np{0,000457}$ est : \quad $4,57 \times 10^{-4}$.

\subsection*{Question 3}
	Comme le triangle est rectangle en B, on choisit [AB] comme base du triangle et donc, la hauteur correspondante est [BC].

	On a donc : \quad $\mathcal{A}_{\mathrm{ABC}} = \dfrac{\text{base} \times \text{hauteur}}{2} = \dfrac{\mathrm{AB} \times \mathrm{BC}}{2} = \dfrac{8 \times 6}{2} = \dfrac{48}{2} = 24$.

	Comme les longueurs sont données en centimètres, le calcul précédent donne une aire en centimètres carrés.

	L'aire du triangle ABC est donc \quad \np[cm^2]{24}.

\subsection*{Question 4}
	Puisque l'on pioche un beignet au hasard et qu'ils semblent identiques, on est en situation d'équiprobabilité.

	La probabilité d'un évènement est le nombre d'issues favorables à l'évènement, divisé par le nombre total d'issues.

	Il y a en tout :\quad $6 + 5 + 4 = 15$ beignets.

	Donc la probabilité de piocher un beignet à la framboise est :\quad $\dfrac{4}{15}$.

\subsection*{Question 5}
	Le coefficient multiplicateur d'une baisse de 10\,\% est : \quad $c = 1 - \dfrac{10}{100} = 1 - 0,1 = 0,9$.

 $800 \times 0,9 = 80 \times 9 =  720$, donc le prix, après réduction, est de: 720\,\euro{}.

	\emph{Variante :} $800 \times \dfrac{10}{100} = 80$, \quad 10\,\% de 800\,\euro{} représentent 80\,\euro{}.

$800 - 80 = 720$, donc le prix, après réduction, est de: 720\,\euro{}.

\subsection*{Question 6}
	Développons : \quad $B = 4y(3y - 1) = 4y \times 3y - 4y \times 1 = 4 \times 3 \times y \times y - 4y = 12y^2 - 4y$.

	La forme développée de $B$ est :\quad $B = 12y^2 - 4y$.

\subsection*{Question 7}
	On sait qu'un litre correspond à \np[dm^3]{1}, et que $\np[dm]{1} = \np[cm]{10}$, donc $\np[dm^3]{1} = 10^3\,\mathrm{cm}^3 = \np[cm^3]{1000}$.


	On peut faire un tableau de conversion :\quad
	\begin{tblr}{colspec={*{6}{X[c,4mm]}},colsep=1pt,hlines, vlines,vline{4}={1.2pt}}
		\SetCell[c=3]{c}{$\mathrm{dm}^3$} &  & & \SetCell[c=3]{c}$\mathrm{cm}^3$ \\
		&&3&5&7&0\\
	\end{tblr}

	\np[L]{3,57} correspond donc à \np[cm^3]{3570}.

	\textbf{Bonne réponse : D}.

\subsection*{Question 8}
	L'image de 4 par la fonction $f$ est :\quad $f(4) = 3 \times 4 - 5 = 12 - 5 = 7$.

	On a $f(4) = 7$.

	\textbf{Bonne réponse : B}.

\subsection*{Question 9}
	ABCD est un quadrilatère.

	Si on appelle O le point d'intersection des diagonales, les codages de la figure indiquent que DO = OB, et AO = OC.

	Ainsi, les diagonales se coupent en leurs milieux : cela garantit que ABCD est un parallélogramme.

	Comme le codage pour AO n'est pas le même que celui pour DO, on ne peut pas affirmer AO = DO : donc les diagonales n'ont pas forcément la même longueur, on ne peut pas affirmer que ABCD est un rectangle.

	Comme il n'y a aucun codage d'angle droit entre les diagonales, on ne peut pas affirmer que ABCD est un losange.

	Comme un carré est à la fois un losange et un rectangle, on ne peut pas affirmer que ABCD est un carré.

	\emph{Variante :} Pour l'absence de codage pour un rectangle, on peut aussi dire qu'il n'y a pas de codage d'angle droit entre deux côtés consécutifs, et pour l'absence de codage pour un losange, on peut aussi dire qu'il n'y a pas de codage de côtés consécutifs de même longueur.

	\medskip

	Quoi qu'il en soit, la seule nature qui est confirmée par les codages, c'est que le quadrilatère ABCD est un parallélogramme.

	\textbf{Bonne réponse : D}.

\needspace{10\baselineskip}
